\documentclass[12pt]{article}

\usepackage{graphicx,url}
\usepackage[T1]{fontenc}
\usepackage[utf8]{inputenc}
\usepackage[brazil]{babel}
\usepackage{sbc-template}
\usepackage{booktabs}
\usepackage{longtable}
\usepackage{tabularx}
\usepackage{array}
\usepackage{placeins}
\usepackage{xcolor}
\usepackage[hidelinks]{hyperref}
\usepackage{pgfplots}
\usepackage{tikz}

\usetikzlibrary{arrows.meta,positioning}
\pgfplotsset{compat=1.18}

\definecolor{GovBlue}{HTML}{005EA8}
\definecolor{GovGold}{HTML}{9A6700}
\definecolor{GovGreen}{HTML}{2E6B45}

\hypersetup{
  pdftitle={Desenvolvimento de uma Aplicação Web para Consulta de Dados de Servidores do Executivo Federal},
  pdfauthor={Pedro Machado, Marcos Freitas e Alexsander Santos}
}

\setlength{\tabcolsep}{3pt}
\renewcommand{\arraystretch}{1.08}
\newcolumntype{Y}{>{\raggedright\arraybackslash}X}
\newcolumntype{P}[1]{>{\raggedright\arraybackslash}p{#1}}
\newcommand{\code}[1]{\texttt{#1}}

\sloppy

\title{Desenvolvimento de uma Aplicação Web para Consulta de Dados de Servidores do Executivo Federal}

\author{Pedro Machado\inst{1}, Marcos Freitas\inst{1},
  Alexsander Santos\inst{1}}

\address{Universidade Federal do Rio Grande (FURG)\\
  Sistemas para Internet II -- Rio Grande -- RS -- Brasil
  \email{pedro.machado@furg.br, mcopellofreitas@hotmail.com}
  \email{alexsander.oml@gmail.com}
}

\begin{document}

\maketitle

\begin{abstract}
  This report describes the development of a full stack application for
  querying public data about active and retired employees of the Brazilian
  Federal Executive Branch. The solution combines PostgreSQL, a Flask API
  served by Gunicorn, and a React interface coordinated by Docker Compose.
  It records daily task allocation, each contributor's Git history, the
  burndown of two sprints, implemented requirements, executed tests, team
  analysis, and verified limitations. The tracked sprint backlog totalled 115
  Story Points; completed tasks are distinguished from assignment requirements
  that remain only partially evidenced.
\end{abstract}

\begin{resumo}
  Este relatório descreve o desenvolvimento de uma aplicação full stack para
  consulta de dados públicos de servidores ativos e aposentados do Poder
  Executivo Federal. A solução combina PostgreSQL, API Flask executada pelo
  Gunicorn e interface React coordenada pelo Docker Compose. São registrados a
  divisão diária das tarefas, o histórico Git de cada integrante, o burndown de
  duas sprints, os requisitos implementados, os testes executados, a análise da
  equipe e as limitações comprovadas. O backlog acompanhado totalizou 115 Story
  Points; as tarefas concluídas são diferenciadas dos requisitos do enunciado
  que permanecem apenas parcialmente comprovados.
\end{resumo}

\section{Introdução e objetivo}

O trabalho propôs uma aplicação full stack para consultar informações de
servidores ativos e aposentados do Poder Executivo Federal. O enunciado exigiu
o uso de GitHub, metodologia ágil, Docker, PostgreSQL, Python, Flask, Gunicorn e
uma interface web, além do registro das tarefas diárias, dos commits, do
burndown e da análise crítica da equipe \cite{enunciado2026}.

As fontes definidas foram os conjuntos Gestão de Pessoas do Executivo Federal
-- Aposentados e Gestão de Pessoas do Executivo Federal -- Carreiras/Cargos
\cite{dadosAposentados,dadosCarreiras}. O escopo obrigatório incluiu consultas
por nome exato, similaridade, cargo, UF, órgão e combinações, com apresentação
de um ou vários resultados. O aplicativo desktop foi mantido fora do escopo
porque corresponde ao item adicional da avaliação.

Este relatório está organizado segundo os itens avaliados: organização da
equipe e sprints; tarefas e histórico diário do GitHub; burndown; arquitetura e
implementação; testes; limitações; pontos fortes e fracos; dificuldades e
impedimentos; e preparação da entrega.

\section{Organização da equipe e metodologia ágil}

\subsection{Responsabilidades}

A equipe foi composta por três integrantes. Pedro Machado exerceu a liderança técnica
e concentrou infraestrutura, garantia de qualidade e documentação de gestão.
Marcos Freitas atuou no banco de dados e no backend. Alexsander Santos ficou responsável pelo
frontend. A Tabela~\ref{tab:equipe} relaciona os nomes usados no relatório às
identidades públicas do histórico Git.

\begin{table}[ht]
\centering
\caption{Divisão de responsabilidades da equipe.}
\label{tab:equipe}
\begin{tabularx}{\textwidth}{P{3.0cm}P{3.7cm}Y}
\toprule
Integrante & Usuário no GitHub & Responsabilidades principais \\
\midrule
Pedro Machado & \code{pedromchd} & Liderança técnica, infraestrutura, QA e documentação de gestão. \\
Marcos Freitas & \code{marcoscfreitas} & Banco de dados, preparação dos dados e backend. \\
Alexsander Santos & \code{alexsander-oml} & Frontend, integração com a API e apresentação dos resultados. \\
\bottomrule
\end{tabularx}
\end{table}

\subsection{Sprints e fluxo de trabalho}

O desenvolvimento foi dividido em duas sprints de 15 dias. A Sprint 1 ocorreu
de 12/09 a 26/09/2026 e a Sprint 2 de 27/09 a 11/10/2026, com 12/10 reservado
para o fechamento anterior à entrega. O Jira registrou descrição, responsável,
Story Points e estado de cada atividade. As tarefas receberam 1 ou 2 pontos e
foram ordenadas pelas dependências entre infraestrutura, dados, backend,
frontend, integração, testes e documentação.

As mudanças foram produzidas em branches identificadas pela área e pela chave
do Jira. Cada commit foi vinculado a uma única tarefa e preservou a autoria do
responsável. Até duas tarefas relacionadas, do mesmo integrante e da mesma
área, puderam compartilhar uma branch, sem mistura no mesmo commit. As
integrações foram realizadas em \code{dev}.

\begin{table}[ht]
\centering
\caption{Distribuição final de Story Points por integrante.}
\label{tab:sp-equipe}
\begin{tabular}{lrrr}
\toprule
Integrante & Sprint 1 & Sprint 2 & Total \\
\midrule
Pedro Machado & 16 & 23 & 39 \\
Marcos Freitas & 22 & 17 & 39 \\
Alexsander Santos & 18 & 19 & 37 \\
\midrule
Total & 56 & 59 & 115 \\
\bottomrule
\end{tabular}
\end{table}

O backlog da Sprint 1 totalizou 56 pontos. O da Sprint 2 começou com 57 pontos
e recebeu 2 pontos em 09/10 para a revisão final do relatório, totalizando 59;
as atividades registradas foram concluídas nessa data. Esse fechamento em
Story Points não substitui a verificação dos requisitos, apresentada na
Tabela~\ref{tab:requisitos}. Dias sem commit foram mantidos no histórico, sem
atribuição retroativa de trabalho.

\section{Tarefas diárias e histórico do GitHub}

As Tabelas~\ref{tab:historico-s1} e~\ref{tab:historico-s2} reúnem, por dia e por
usuário, as tarefas delegadas e os commits correspondentes. Assim, os IDs do
backlog não aparecem isolados: cada registro informa também o resultado
produzido. O histórico foi conferido em \code{dev} até \code{033a83b}, primeira
consolidação da PROJ-72. A presente correção continua vinculada à mesma tarefa;
seu próprio hash não pode constar no conteúdo que o determina.

\small
\begin{longtable}{P{1.35cm}P{3.8cm}P{9.15cm}}
\caption{Tarefas e commits diários da Sprint 1.}\label{tab:historico-s1}\\
\toprule
Data & Integrante & Tarefas delegadas e commits \\
\midrule
\endfirsthead
\toprule
Data & Integrante & Tarefas delegadas e commits \\
\midrule
\endhead
12/09 & Pedro Machado & \code{INF-01/PROJ-3}: configurar e ajustar a estrutura inicial do repositório -- \code{82574ee}, \code{0a50fb5}. \\
12/09 & Marcos Freitas e Alexsander Santos & Sem entrega integrada. \\
13/09 & Pedro Machado & \code{INF-02/PROJ-4}: organizar backlog e quadro do Jira -- atividade de gestão sem commit. \\
13/09 & Marcos Freitas & \code{BD-01/PROJ-15}: adicionar \code{psycopg[binary]} às dependências do backend para conexão com PostgreSQL -- \code{1148065}; \code{BD-02/PROJ-16}: analisar os datasets -- \code{882f94d}. \\
13/09 & Alexsander Santos & \code{FRONT-01/PROJ-29}: configurar a estrutura do frontend -- \code{58ffbc9}; \code{FRONT-02/PROJ-30}: definir a identidade visual -- \code{776d42f}. \\
14/09 & Pedro Machado & \code{INF-03/PROJ-5}: criar a composição inicial -- \code{c8395e6}. \\
14/09 & Marcos Freitas & \code{BD-03/PROJ-17}: mapear os campos necessários -- \code{d64f588}. \\
14/09 & Alexsander Santos & \code{FRONT-03/PROJ-31}: criar o layout base -- \code{9ab57fe}. \\
15/09 & Pedro Machado & \code{INF-04/PROJ-6}: criar a imagem do banco -- \code{9785ee7}. \\
15/09 & Marcos Freitas & \code{BD-04/PROJ-18}: modelar o schema -- \code{fe3b0d8}. \\
15/09 & Alexsander Santos & \code{FRONT-04/PROJ-32}: criar o formulário de busca -- \code{ea59923}. \\
16/09 & Pedro Machado & \code{INF-05/PROJ-7}: criar a imagem do backend -- \code{e4be7c3}. \\
16/09 & Marcos Freitas & \code{BD-05/PROJ-19}: criar as tabelas do banco -- \code{286d627}. \\
16/09 & Alexsander Santos & \code{FRONT-05/PROJ-33}: estilizar o formulário -- \code{9febd91}. \\
17/09 & Equipe & Sem entrega integrada. \\
18/09 & Pedro Machado & \code{INF-06/PROJ-8}: criar a imagem do frontend -- \code{580e9ef}; \code{INF-07/PROJ-9}: validar a orquestração -- \code{3baff68}. \\
18/09 & Marcos Freitas & \code{BD-06/PROJ-20}: preparar os dados para importação -- \code{70c9196}. \\
18/09 & Alexsander Santos & \code{FRONT-06/PROJ-34}: criar a listagem fictícia -- \code{4123f17}; \code{FRONT-07/PROJ-35}: estilizar a listagem -- \code{f6a31fd}. \\
19/09 & Equipe & Sem entrega integrada. \\
20/09 & Pedro Machado & \code{INF-08/PROJ-10}: configurar variáveis de ambiente -- \code{493c300}; \code{QA-01/PROJ-11}: definir o plano de testes -- \code{edf7e9d}. \\
20/09 & Marcos Freitas & \code{BD-07/PROJ-21}: importar dados tratados -- \code{d27e820}; \code{BD-08/PROJ-22}: criar índices -- \code{66f5f0b}; \code{BACK-01/PROJ-23}: criar o esqueleto Flask -- \code{e4513c3}. \\
20/09 & Alexsander Santos & \code{FRONT-08/PROJ-36}: criar a paginação fictícia -- \code{8570123}; \code{FRONT-09/PROJ-37}: criar estados simulados -- \code{ad846df}. \\
21--22/09 & Equipe & Sem entrega integrada. \\
23/09 & Pedro Machado & \code{QA-02/PROJ-12}: definir casos de SQL injection -- \code{b12a706}; \code{LID-01/PROJ-13}: registrar o acompanhamento -- \code{79ef587}; \code{LID-02/PROJ-14}: estruturar o relatório -- \code{26cd934}. \\
23/09 & Marcos Freitas & \code{BACK-02/PROJ-24}: configurar Gunicorn multithread -- \code{ccdf5c4}; \code{BACK-03/PROJ-25}: criar conexão com PostgreSQL -- \code{4cfa188}; \code{BACK-04/PROJ-26}: implementar busca exata -- \code{70f63a5}. \\
23/09 & Alexsander Santos & \code{FRONT-10/PROJ-38}: ajustar a responsividade -- \code{44d97f8}. \\
24--25/09 & Equipe & Sem entrega integrada. \\
26/09 & Pedro Machado & \code{LID-01/PROJ-13}: encerrar o acompanhamento da Sprint 1 -- \code{5236647}. \\
26/09 & Marcos Freitas & \code{BACK-05/PROJ-27}: implementar filtros por cargo e UF -- \code{8e440b9}; \code{BACK-06/PROJ-28}: implementar filtro por órgão -- \code{47e3872}. \\
26/09 & Alexsander Santos & Sem entrega integrada. \\
\bottomrule
\end{longtable}
\normalsize

\small
\begin{longtable}{P{1.35cm}P{3.8cm}P{9.15cm}}
\caption{Tarefas e commits diários da Sprint 2.}\label{tab:historico-s2}\\
\toprule
Data & Integrante & Tarefas delegadas e commits \\
\midrule
\endfirsthead
\toprule
Data & Integrante & Tarefas delegadas e commits \\
\midrule
\endhead
27/09 & Pedro Machado & \code{INF-09/PROJ-39}: integrar a aplicação ao Compose -- \code{467b7a3}. \\
27/09 & Marcos Freitas & \code{BACK-07/PROJ-51}: implementar busca combinada -- \code{f8a1064}. \\
27/09 & Alexsander Santos & \code{FRONT-11/PROJ-61}: integrar o frontend à API -- \code{cc9f180}; \code{FRONT-12/PROJ-62}: implementar busca exata -- \code{3703715}. \\
28--29/09 & Equipe & Sem entrega integrada. \\
30/09 & Pedro Machado & \code{INF-10/PROJ-40}: configurar healthchecks e reinício -- \code{fcc9fb3}; \code{INF-11/PROJ-41}: limitar buscas simultâneas -- \code{555a6a3}; \code{INF-12/PROJ-42}: aplicar timeout -- \code{dd80b71}. \\
30/09 & Marcos Freitas & \code{BACK-08/PROJ-52}: implementar similaridade -- \code{ed9e38f}; \code{BACK-09/PROJ-53}: implementar paginação -- \code{389c272}; \code{BACK-10/PROJ-54}: validar parâmetros -- \code{728568f}. \\
30/09 & Alexsander Santos & \code{FRONT-13/PROJ-63}: implementar similaridade -- \code{ef2df83}; \code{FRONT-14/PROJ-64}: implementar filtros individuais -- \code{4b13c89}; \code{FRONT-15/PROJ-65}: implementar busca combinada -- \code{fb65451}. \\
01--03/10 & Equipe & Sem entrega integrada. \\
04/10 & Pedro Machado & \code{QA-03/PROJ-43}: executar carga -- \code{1274a0d}; \code{QA-04/PROJ-44}: registrar resultados e limitações -- \code{569da1f}; \code{QA-05/PROJ-45}: executar cenários de erro -- \code{98f891d}; \code{QA-06/PROJ-46}: validar SQL injection -- \code{49fe3e6}. \\
04/10 & Marcos Freitas & \code{BACK-11/PROJ-55}: tratar erros -- \code{545c692}; \code{BACK-12/PROJ-56}: padronizar JSON -- \code{e1d822f}; \code{BD-09/PROJ-57}: otimizar similaridade -- \code{1404a59}; \code{BACK-13/PROJ-58}: adicionar testes -- \code{41eb8da}. \\
04/10 & Alexsander Santos & \code{FRONT-16/PROJ-66}: integrar paginação -- \code{a9ed8db}; \code{FRONT-17/PROJ-67}: diferenciar resultado único -- \code{09b01a6}; \code{FRONT-18/PROJ-68}: adicionar feedback -- \code{982d69a}; \code{FRONT-19/PROJ-69}: ajustar responsividade -- \code{70e40cd}. \\
05/10 & Equipe & Sem entrega integrada. \\
06/10 & Pedro Machado & \code{LID-03/PROJ-47}: consolidar tarefas e histórico -- \code{3acd6e5}; \code{LID-04/PROJ-48}: criar o burndown -- \code{f2faee7}. \\
06/10 & Marcos Freitas & \code{BACK-14/PROJ-59}: documentar a API -- \code{24a866c}; \code{BACK-15/PROJ-60}: corrigir e revalidar respostas -- \code{0a28f98}. \\
06/10 & Alexsander Santos & \code{FRONT-20/PROJ-70}: corrigir estados de busca e respostas inválidas -- \code{6ab7768}. \\
07--08/10 & Equipe & Sem entrega integrada; Marcos Freitas e Alexsander Santos já não possuíam pendências obrigatórias. \\
09/10 & Pedro Machado & \code{LID-05/PROJ-49}: registrar a análise crítica -- \code{cc1ebfa}; \code{LID-06/PROJ-50}: permitir a inicialização sem arquivo de ambiente e validar a aplicação em clone limpo -- \code{e5647e8}; \code{LID-07/PROJ-72}: consolidar o relatório -- \code{033a83b}. \\
09/10 & Marcos Freitas e Alexsander Santos & Sem pendências obrigatórias. \\
10/10 & Pedro Machado & \code{LID-07/PROJ-72}: revisar o relatório com o pacote oficial da SBC e alinhá-lo ao enunciado -- presente revisão. \\
10/10 & Marcos Freitas e Alexsander Santos & Sem pendências obrigatórias. \\
11/10 & Equipe & Período reservado para fechamento. \\
\bottomrule
\end{longtable}
\normalsize

\subsection{Burndown}

O saldo real representa o fim de cada dia. A Sprint 2 mantém a linha-base de
57 pontos e registra a alteração de escopo de 09/10 separadamente: havia 3
pontos restantes, a PROJ-72 acrescentou 2 e as três tarefas do dia concluíram
5 pontos, levando o saldo a zero.

\begin{figure}[ht]
\centering
\begin{tikzpicture}
\begin{axis}[
  width=0.47\textwidth,
  height=6.0cm,
  title={Sprint 1},
  xmin=0,xmax=15,ymin=0,ymax=60,
  xtick={0,3,6,9,12,15},
  xticklabels={início,14/09,17/09,20/09,23/09,26/09},
  ytick={0,15,30,45,60},
  grid=major,
  tick label style={font=\scriptsize},
  title style={font=\small\bfseries},
  legend style={font=\scriptsize,at={(0.5,-0.25)},anchor=north,legend columns=2},
  xlabel={dia},ylabel={SP restantes}
]
\addplot[GovGold,dashed,thick] coordinates {(0,56) (15,0)};
\addplot[GovBlue,thick,mark=*] coordinates {
  (0,56) (1,55) (2,48) (3,43) (4,38) (5,34) (6,34) (7,25)
  (8,25) (9,14) (10,14) (11,14) (12,4) (13,4) (14,4) (15,0)
};
\legend{ideal,real}
\end{axis}
\end{tikzpicture}\hfill
\begin{tikzpicture}
\begin{axis}[
  width=0.47\textwidth,
  height=6.0cm,
  title={Sprint 2},
  xmin=0,xmax=15,ymin=0,ymax=60,
  xtick={0,3,6,9,12,15},
  xticklabels={início,29/09,02/10,05/10,08/10,11/10},
  ytick={0,15,30,45,60},
  grid=major,
  tick label style={font=\scriptsize},
  title style={font=\small\bfseries},
  xlabel={dia},ylabel={SP restantes}
]
\addplot[GovGold,dashed,thick] coordinates {(0,57) (15,0)};
\addplot[GovBlue,thick,mark=*] coordinates {
  (0,57) (1,49) (2,49) (3,49) (4,32) (5,32) (6,32) (7,32)
  (8,11) (9,11) (10,3) (11,3) (12,3) (13,0)
};
\node[anchor=south,font=\scriptsize,align=center,text=GovGreen]
  at (axis cs:10.7,13) {+2 SP no escopo\\5 SP concluídos};
\end{axis}
\end{tikzpicture}
\caption{Story Points restantes nas duas sprints. A Sprint 2 terminou com escopo final de 59 SP.}
\label{fig:burndown}
\end{figure}

\begin{table}[ht]
\centering
\caption{Fechamento das sprints.}
\label{tab:fechamento}
\begin{tabular}{lrrrrl}
\toprule
Sprint & Base & Ajuste & Escopo final & Concluídos & Fechamento \\
\midrule
Sprint 1 & 56 & 0 & 56 & 56 & 26/09 \\
Sprint 2 & 57 & +2 & 59 & 59 & 09/10 \\
\bottomrule
\end{tabular}
\end{table}
\FloatBarrier

\section{Arquitetura e implementação}

\subsection{Docker e execução}

A aplicação possui três serviços orquestrados pelo Docker Compose:

\begin{itemize}
  \item \code{db}: PostgreSQL 18 Alpine, responsável pelo schema, índices e
  extensão \code{pg\_trgm};
  \item \code{backend}: Python 3.14, Flask 3.1.3, Gunicorn 26.2.0 e psycopg
  3.3.5;
  \item \code{frontend}: React 19, TypeScript, Vite 8 e Tailwind CSS 4.
\end{itemize}

\begin{figure}[ht]
\centering
\begin{tikzpicture}[
  node distance=0.52cm,
  box/.style={draw=GovBlue, rounded corners=2pt, fill=blue!4,
    minimum height=1.30cm, text width=2.65cm, align=center},
  arrow/.style={-{Latex[length=2.2mm]}, thick, draw=GovBlue}
]
\node[box] (browser) {Navegador\\porta 5173};
\node[box, right=of browser] (front) {Frontend\\React e Vite};
\node[box, right=of front] (back) {Backend\\Flask e Gunicorn};
\node[box, right=of back] (db) {Banco\\PostgreSQL};
\draw[arrow] (browser) -- node[above, font=\scriptsize] {HTTP} (front);
\draw[arrow] (front) -- node[above, font=\scriptsize] {\code{/api}} (back);
\draw[arrow] (back) -- node[above, font=\scriptsize] {SQL} (db);
\end{tikzpicture}
\caption{Comunicação entre os módulos da aplicação.}
\label{fig:arquitetura}
\end{figure}

O navegador acessa o frontend pela porta 5173. As requisições iniciadas em
\code{/api} são encaminhadas pelo proxy do Vite ao backend na porta interna
8000. O backend acessa o PostgreSQL pelo nome de serviço \code{db}; o banco não
publica porta para o host. O Gunicorn usa um processo e quatro threads. O
Compose limita a API a duas buscas simultâneas e aplica timeout de 2.000 ms.
Banco e backend possuem healthchecks e os três serviços usam
\code{restart: unless-stopped}.

A execução a partir da raiz usa o comando recomendado pela documentação do
Docker Compose \cite{dockerCompose}:

\begin{center}
\code{docker compose up -{}-build}
\end{center}

O comando constrói e inicia os três serviços. Entretanto, um clone limpo cria
o schema sem registros públicos. Portanto, a infraestrutura fica disponível,
mas consultas úteis dependem da execução prévia do tratamento e da importação.
Essa diferença em relação à expressão ``pronta para uso'' do enunciado é uma
limitação da entrega, não um resultado considerado aprovado.

\subsection{Tratamento dos dados e PostgreSQL}

Foram analisados os conjuntos de Aposentados e Carreiras/Cargos e seus
dicionários de dados \cite{dicionarioAposentados,dicionarioCarreiras}. A
inspeção abrangeu as competências de setembro de 2025 a agosto de 2026. Os
recursos de agosto de 2026 estavam indisponíveis e o arquivo de Aposentados de
dezembro de 2025 estava vazio.

Os campos utilizados foram nome, cargo, órgão e UF. CPF, matrícula,
remuneração e os demais dados pessoais ou funcionais não foram incorporados. A
UF não existe no conjunto de Aposentados e permanece nula nesses registros. Os
arquivos usam Latin-1 e separação por ponto e vírgula.

O tratamento remove bytes NUL e espaços excedentes; ignora linhas vazias,
cabeçalhos e linhas de identificação; trata o delimitador final adicional de
Carreiras/Cargos; valida quantidade de colunas e campos obrigatórios;
normaliza a UF; rejeita UFs malformadas; remove registros idênticos no mesmo
arquivo; e gera CSV UTF-8 com origem, arquivo e linha do registro.

O modelo relacional contém \code{servidor}, \code{cargo}, \code{orgao} e
\code{uf}. A importação usa tabela temporária, insere os domínios distintos e
cria os relacionamentos em \code{servidor}. O processo exige tabelas de destino
vazias e interrompe a importação se as quantidades divergirem. Não existe
evidência persistida da contagem final de uma importação pública; por isso, este
relatório não atribui um total importado.

A extensão \code{pg\_trgm} sustenta a similaridade. O nome do servidor usa
índice GiST, compatível com ordenação por distância textual. Cargo e órgão
possuem índices trigram GIN, e as chaves estrangeiras usam índices B-tree.

\subsection{Backend e API}

A API disponibiliza \code{GET /health} e \code{GET /api/servidores}. A consulta
aceita filtros exatos por nome, cargo, UF e órgão, similaridade por nome e
combinações entre filtros. No parâmetro \code{similar}, os marcadores
\code{\%} dos formatos exemplificados no enunciado são removidos antes da
busca. O limite de similaridade varia de 1 a 100, com padrão 100. A paginação
usa \code{page} e \code{page\_size}, com até 100 registros por página.

As respostas usam um envelope JSON com \code{sucesso}, \code{mensagem} e
\code{dados}; respostas paginadas incluem \code{paginacao}. A API diferencia
registro único, múltiplos registros, ausência de resultados e página
inexistente. UF, números e parâmetros desconhecidos são validados antes do
acesso ao banco.

As consultas separam parâmetros do SQL e entradas compatíveis com injeção são
recusadas. A API trata entrada inválida, recurso inexistente, método não
permitido, excesso de concorrência, falha interna e timeout. A terceira busca
simultânea recebe HTTP 429 e o cancelamento pelo PostgreSQL é convertido em
HTTP 504.

\subsection{Frontend}

O frontend apresenta nome, modo exato ou similar, limite da similaridade,
cargo, UF e órgão. Os filtros podem ser usados isoladamente ou em conjunto. A
interface não consulta ao abrir e impede envio sem filtro.

Cada consulta usa páginas de dez registros. A navegação oferece anterior,
próxima e até cinco números de página. A interface implementa estados para
resultado único, lista, carregamento, ausência de resultados e erro. Respostas
inválidas e falhas de rede produzem mensagens controladas. O layout é
responsivo, usa rótulos, regiões de status e alertas e toma o Portal Brasileiro
de Dados Abertos como referência visual.

O código desses estados está implementado e o build foi validado. A evidência
integrada preservada, contudo, cobre apenas o caminho de uma busca sem
resultados em banco vazio. Não há evidência persistida de uma bateria manual do
frontend com dados positivos, resultado único, múltiplos resultados e
paginação; essa lacuna é registrada na Seção~\ref{sec:limitacoes}.

\section{Atendimento aos requisitos do enunciado}

A Tabela~\ref{tab:requisitos} separa implementação, evidência e limitação. O
estado ``parcial'' não invalida o código existente; indica que a evidência
preservada não comprova integralmente o comportamento exigido.

\small
\begin{longtable}{P{4.1cm}P{7.4cm}P{2.1cm}}
\caption{Relação entre requisitos e evidências.}\label{tab:requisitos}\\
\toprule
Requisito & Implementação ou evidência & Situação \\
\midrule
\endfirsthead
\toprule
Requisito & Implementação ou evidência & Situação \\
\midrule
\endhead
GitHub, equipe e metodologia ágil & Três identidades, tasks Jira, commits atômicos, duas sprints e histórico diário. & Atendido \\
Três containers e Dockerfiles & Banco, backend e frontend construídos e iniciados pelo Compose. & Atendido \\
Aplicação pronta após \code{up -{}-build} & Serviços iniciam, mas o banco público exige tratamento e importação posteriores. & Parcial \\
PostgreSQL, preparação e índices & Schema, pipeline, validações e índices existem; falta contagem persistida da importação pública final. & Parcial \\
Python, Flask e Gunicorn & API executada pelo Gunicorn e validada com PostgreSQL real. & Atendido \\
Filtros, combinações, JSON, paginação e erros & Implementação e baterias automatizada, funcional, de timeout, concorrência e segurança. & Atendido \\
Frontend apresentando os dados & Build e caminho integrado sem resultados foram verificados; cenários positivos não possuem evidência persistida. & Parcial \\
Sobrecarga e proteção & Limite de concorrência, HTTP 429, timeout e disponibilidade posterior foram exercitados. & Atendido \\
Tarefas, commits, burndown e análise da equipe & Tabelas diárias, Figura~\ref{fig:burndown} e Seção~\ref{sec:analise}. & Atendido \\
Aplicativo desktop & Item adicional não desenvolvido. & Fora do escopo obrigatório \\
\bottomrule
\end{longtable}
\normalsize

\section{Estratégia e resultados dos testes}

\begin{table}[ht]
\centering
\caption{Consolidação das verificações executadas.}
\label{tab:testes}
\small
\begin{tabularx}{\textwidth}{P{4.0cm}Y}
\toprule
Verificação & Resultado \\
\midrule
Backend automatizado & 12 de 12 testes aprovados. \\
API funcional e erros & 20 de 20 cenários aprovados contra Flask, Gunicorn e PostgreSQL reais. \\
SQL injection & 45 de 45 casos aprovados, sem alteração de tabelas, contagens ou fingerprints. \\
Carga de capacidade & 20 requisições, concorrência 2, todas HTTP 200 e sem erro de transporte. \\
Sobrecarga & 200 requisições, concorrência 32; 2 HTTP 200 e 198 HTTP 429, sem erro de transporte. \\
Disponibilidade posterior & \code{/health} respondeu HTTP 200 após a sobrecarga. \\
Timeout real & Consulta bloqueada retornou HTTP 504 em aproximadamente 2 segundos. \\
Clone limpo & Três imagens construídas; banco e backend saudáveis; frontend e \code{/health} com HTTP 200. \\
Integração do frontend & Busca percorreu frontend, proxy e API e apresentou o estado sem resultados. \\
\bottomrule
\end{tabularx}
\end{table}

A bateria funcional cobriu filtro vazio, espaços, tentativa representativa de
injeção, nome inexistente, múltiplos resultados paginados, cargo, UF e órgão
inválidos, rota e página inexistentes, parâmetro desconhecido, limites
inválidos, combinação inválida, timeout, concorrência e disponibilidade
posterior. Na bateria de segurança, os payloads foram enviados aos cinco
filtros; as tabelas, contagens e fingerprints permaneceram inalteradas.

No cenário de capacidade, 20 respostas HTTP 200 foram concluídas em 2,921 s,
com 6,85 requisições por segundo, mediana de 276,08 ms e p95 de 367,76 ms. Na
sobrecarga, o limite recusou 198 requisições com HTTP 429 e permitiu duas
respostas HTTP 200. Não houve erro de transporte e a saúde permaneceu
disponível.

\subsection{Limitações}\label{sec:limitacoes}

As limitações comprovadas são:

\begin{itemize}
  \item o clone limpo inicia os serviços e cria o schema, mas não carrega os
  arquivos públicos automaticamente;
  \item não existe contagem persistida da importação pública final;
  \item a evidência integrada do frontend cobre banco vazio, não os cenários
  positivos completos;
  \item a execução final em clone limpo de 09/10 está resumida neste relatório,
  mas não possui log independente versionado; \path{docs/validacao-compose.md}
  registra somente a validação preliminar de 18/09;
  \item as baterias completas funcional, de injeção e de carga foram executadas
  em 04/10, antes da correção do backend de 06/10; depois dela, os 12 testes
  automatizados passaram em 09/10, mas as baterias completas não foram
  repetidas;
  \item os testes ocorreram localmente, com uma instância do backend e, nos
  testes de carga, 100.000 registros sintéticos;
  \item não houve teste prolongado de estabilidade nem distribuição entre
  múltiplas instâncias;
  \item a capacidade configurada é de duas buscas simultâneas e rejeita o
  excedente com HTTP 429;
  \item o aplicativo desktop adicional não foi desenvolvido.
\end{itemize}

\section{Análise crítica da equipe}\label{sec:analise}

\subsection{Pontos fortes}

\begin{itemize}
  \item A divisão por área permaneceu estável e o histórico preservou a autoria
  dos três integrantes.
  \item As tarefas permaneceram identificadas por Jira, mesmo quando atividades
  relacionadas compartilharam a mesma integração.
  \item As evidências reproduzíveis cobriram API, segurança, carga e clone
  limpo. Problemas encontrados pelo QA produziram correções no backend e no
  frontend.
\end{itemize}

\subsection{Pontos fracos}

\begin{itemize}
  \item A cadência não foi uniforme e houve quedas concentradas no burndown.
  \item A capacidade de duas buscas simultâneas protege a aplicação, mas é baixa
  para picos de acesso.
  \item A documentação final ficou concentrada no encerramento, exigindo
  conferência simultânea do Jira, histórico Git e evidências.
  \item A importação pública e os cenários positivos do frontend não ficaram
  registrados com evidências finais reproduzíveis.
\end{itemize}

\subsection{Dificuldades e impedimentos}

\small
\begin{tabularx}{\textwidth}{P{2.0cm}YY}
\toprule
Período & Dificuldade ou impedimento & Ação e resultado \\
\midrule
12/09 & Estrutura inicial diferente da convenção do enunciado. & Backend, frontend, banco e composição foram mantidos diretamente na raiz do repositório. \\
13/09 & Três integrantes no mesmo computador precisavam preservar identidades distintas. & Clones e identidades separadas mantiveram a autoria dos commits. \\
15/09 & A validação das imagens dependia do Docker em execução. & O ambiente passou a ser iniciado antes das verificações. \\
17--25/09 & Interrupções e dias sem integração acumularam entregas. & O trabalho foi retomado em lotes, com commits atômicos e chaves Jira. \\
04/10 & A carga confirmou baixa capacidade e custo da busca ampla por similaridade. & Limite, timeout e índice fizeram o excesso falhar de maneira controlada. \\
10/10 & O relatório inicial não aplicava o pacote oficial da SBC nem priorizava os itens avaliados. & O documento foi recomposto com o template oficial e reorganizado pelo enunciado. \\
\bottomrule
\end{tabularx}
\normalsize

\section{Preparação da entrega}

O repositório mantém \code{backend/}, \code{frontend/}, \code{db/},
\code{docs/} e \code{docker-compose.yml} diretamente na raiz. Para atender à
estrutura do arquivo final sem duplicar diretórios durante o desenvolvimento,
o ZIP de entrega deve envolver essa raiz em uma única pasta externa chamada
\code{projeto/}. Os arquivos dos conjuntos públicos e o banco preenchido não
devem ser incluídos.

Depois da extração, o avaliador deverá entrar em \code{projeto/} e executar
\code{docker compose up -{}-build}. O procedimento inicia os serviços. Como
registrado nas limitações, a realização de buscas com resultados depende da
preparação e importação dos dados públicos.

\section{Conclusão}

O projeto produziu uma aplicação web modularizada em PostgreSQL, API Flask com
Gunicorn e frontend React. Foram implementadas buscas por nome exato,
similaridade, cargo, UF, órgão e combinações, além de paginação, respostas JSON
padronizadas, validação, limite de concorrência, timeout e tratamento dos erros
previstos.

A Sprint 1 encerrou com 56 de 56 pontos e a Sprint 2 com 59 de 59, totalizando
115 Story Points do backlog registrados como concluídos. Esse fechamento do
planejamento não significa atendimento integral aos requisitos classificados
como parciais na Tabela~\ref{tab:requisitos}. As tarefas e os commits dos três
integrantes estão documentados por dia, juntamente com o burndown e a análise
crítica.

As evidências de 04/10 confirmam o backend, a proteção contra injeção e o
controle de sobrecarga naquela versão; os 12 testes automatizados passaram
novamente na validação de clone limpo de 09/10. As baterias completas não foram
repetidas após a correção de 06/10. A avaliação final deve considerar também a
ausência de carga automática dos dados públicos, de contagem final importada e
de evidência positiva completa do frontend. O aplicativo desktop permaneceu
fora do escopo obrigatório.

\section{Disponibilidade do código e evidências}

O código-fonte e o histórico Git estão disponíveis em
\url{https://github.com/marcoscfreitas/spi-ii-3bim}. As evidências usadas neste
relatório incluem:

\begin{itemize}
  \item \path{docs/analise-datasets.md}, \path{docs/mapeamento-campos.md} e
  \path{docs/modelo-banco.md};
  \item \path{backend/README.md}, \path{docs/plano-testes.md} e
  \path{docs/casos-sql-injection.md};
  \item a validação preliminar da composição, de 18/09, em
  \path{docs/validacao-compose.md}; a validação final de clone limpo, de 09/10,
  está resumida neste relatório e no histórico da PROJ-50, sem log independente;
  \item o acompanhamento em \path{docs/acompanhamento-sprint-1.md} e o
  gráfico em \path{docs/burndown.svg};
  \item os registros JSON em \path{docs/evidencias/}, produzidos em 04/10/2026;
  \item o histórico integrado da branch \code{dev}, usado para conferir autoria,
  datas e mensagens.
\end{itemize}

\bibliographystyle{sbc}
\bibliography{relatorio-final}

\end{document}
